%-*-latex-*-
%-*-latex-*-
\newcommand\COURSE{ciss350}
\newcommand\ASSESSMENT{q10201}
\newcommand\ASSESSMENTTYPE{Quiz}
\newcommand\POINTS{\textwhite{xxx/xxx}}

\input{myquizpreamble}
\input{yliow}
\input{\COURSE}
%-*-latex-*-

%-*-latex-*-

%-*-latex-*-
\renewcommand\TITLE{\ASSESSMENTTYPE \ \ASSESSMENT}

\renewcommand\EMAIL{}
\renewcommand\AUTHOR{eagraham1@cougars.ccis.edu}

\textwidth=6in
\begin{document}
\topmatterthree

\LATEXHELPTHREEFIVEZERO

\TURNPAGE

%------------------------------------------------------------------------------
\newpage
\textsc{Runtime for array assignment}

The goal is to compute the big-O runtime of the basic array
assignment algorithm (i.e., copy values from one array to another).

For the next few questions,
the goal is to analyze the runtime of this algorithm:
\begin{Verbatim}[frame=single]
for i = 0, ..., n - 1:
    x[i] = y[i]
\end{Verbatim}

%------------------------------------------------------------------------------
\nextq
Rewrite the above algorithm with
goto statements, assign constant times to each statement and
compute the time taken as a function of $n$, $t_1$, $t_2$, $t_3$, ....
where $t_1, t_2, t_3, ...$ are runtimes for each statement in your
algorithm. Make sure you
number these times starting with $t_1$, then
$t_2$, etc.
Include the number of times each statement is executed.
\\
\ANSWER
\begin{answercode}
          i = 0                   t1      1
LOOP:     if i >= n - 1:          t2      n
              goto ENDLOOP        t3      1
\end{answercode}

%------------------------------------------------------------------------------
\nextq
Let $T(n)$ be the runtime function of the above algorithm.
Write $T(n)$ as a polynomial in descending power of $n$, i.e.,
the \textit{highest degree term appearing first}.
\\
\ANSWER
\begin{answerlong}
$T(n) = (t_{2})n + (t_{1} + t_{3})$
\end{answerlong}
(Correct the above.)

%------------------------------------------------------------------------------
\nextq
Simplify the runtime function $T(n)$
by giving names $A$, $B$, $C$, $D$, ...
to the constants of the function from (a). (The $t_1, t_2, ...$ should
disappear.)
\\
\ANSWER
\begin{answerlong}
$T(n) = A n + B n$
\end{answerlong}
(Correct the above.)

%------------------------------------------------------------------------------
\nextq
Fudge away the constants and write down the simplest $g(n)$ such
that the runtime is a big-O of your $g(n)$.
Your $g(n)$ should be either $n$ or $n^2$ or $n^3$ or $n^{102.25}$ ...
\\
\ANSWER
\begin{answerlong}
$T(n) = O(n)$
\end{answerlong}
(Correct the above.)
\qed

%------------------------------------------------------------------------------
%-*-latex-*-
\newpage
\input{instructions}
\end{document}

